\chapter{Análise de Requisitos}
\label{chap:requisitos}

\section{Introdução}
\label{sec:req-intro}

Este capítulo apresenta a análise de requisitos do sistema, incluindo os requisitos funcionais (RF) e não-funcionais (RNF), organizados por módulo. A priorização segue a metodologia MoSCoW (\textit{Must, Should, Could, Won't}), permitindo uma gestão eficaz do âmbito do projeto ao longo dos sprints de desenvolvimento.


\section{Requisitos Funcionais}
\label{sec:req-funcionais}

Os requisitos funcionais estão organizados em sete módulos, correspondentes às principais áreas funcionais do sistema.

\subsection{Módulo 1 --- Autenticação e Gestão de Utilizadores}

\begin{table}[H]
\centering
\caption{Requisitos funcionais do módulo de autenticação.}
\label{tab:rf-auth}
\small
\begin{tabularx}{\textwidth}{l X c c}
\toprule
\textbf{ID} & \textbf{Descrição} & \textbf{Prioridade} & \textbf{Sprint} \\
\midrule
RF01 & O sistema deve permitir registo com email e password & Must & 3 \\
RF02 & O sistema deve suportar login via OAuth2 (Google) & Must & 3 \\
RF03 & O sistema deve implementar autenticação \ac{2FA} (TOTP) & Should & 3 \\
RF04 & O utilizador deve poder editar o seu perfil (nome, avatar, moeda) & Must & 3 \\
RF05 & O sistema deve permitir recuperação de password por email & Must & 3 \\
\bottomrule
\end{tabularx}
\end{table}

\subsection{Módulo 2 --- Open Banking e Contas}

\begin{table}[H]
\centering
\caption{Requisitos funcionais do módulo de Open Banking.}
\label{tab:rf-openbanking}
\small
\begin{tabularx}{\textwidth}{l X c c}
\toprule
\textbf{ID} & \textbf{Descrição} & \textbf{Prioridade} & \textbf{Sprint} \\
\midrule
RF06 & O sistema deve ligar contas bancárias via Salt Edge \ac{API} & Must & 4 \\
RF07 & O sistema deve importar transações automaticamente (até 12 meses) & Must & 4 \\
RF08 & O sistema deve mostrar saldos atualizados de todas as contas & Must & 4 \\
RF09 & O utilizador deve poder desligar uma conta a qualquer momento & Must & 4 \\
RF10 & O sistema deve sincronizar transações periodicamente (mín. 1x/dia) & Should & 5 \\
RF11 & O sistema deve suportar pelo menos 3 bancos PT (CGD, BCP, Santander) & Must & 4 \\
\bottomrule
\end{tabularx}
\end{table}

\subsection{Módulo 3 --- Dashboard e Visualização}

\begin{table}[H]
\centering
\caption{Requisitos funcionais do módulo de dashboard.}
\label{tab:rf-dashboard}
\small
\begin{tabularx}{\textwidth}{l X c c}
\toprule
\textbf{ID} & \textbf{Descrição} & \textbf{Prioridade} & \textbf{Sprint} \\
\midrule
RF12 & O dashboard deve mostrar saldo total agregado de todas as contas & Must & 5 \\
RF13 & O dashboard deve exibir gráfico de gastos por categoria & Must & 5 \\
RF14 & O dashboard deve mostrar evolução de saldo ao longo do tempo & Must & 5 \\
RF15 & O dashboard deve listar as últimas transações com pesquisa e filtros & Must & 5 \\
RF16 & O sistema deve gerar um relatório mensal automático & Should & 8 \\
RF17 & O utilizador deve poder exportar dados em CSV & Could & 8 \\
\bottomrule
\end{tabularx}
\end{table}

\subsection{Módulo 4 --- Categorização com IA}

\begin{table}[H]
\centering
\caption{Requisitos funcionais do módulo de categorização com IA.}
\label{tab:rf-ia}
\small
\begin{tabularx}{\textwidth}{l X c c}
\toprule
\textbf{ID} & \textbf{Descrição} & \textbf{Prioridade} & \textbf{Sprint} \\
\midrule
RF18 & O sistema deve categorizar transações automaticamente com \ac{ML} & Must & 6 \\
RF19 & O utilizador deve poder corrigir categorias manualmente & Must & 6 \\
RF20 & O modelo \ac{ML} deve aprender com as correções do utilizador & Should & 6 \\
RF21 & O sistema deve ter categorias pré-definidas para o mercado PT & Must & 6 \\
RF22 & O sistema deve detetar transações recorrentes (subscrições, salário) & Should & 6 \\
\bottomrule
\end{tabularx}
\end{table}

\subsection{Módulo 5 --- Orçamentos e Metas}

\begin{table}[H]
\centering
\caption{Requisitos funcionais do módulo de orçamentos e metas.}
\label{tab:rf-budgets}
\small
\begin{tabularx}{\textwidth}{l X c c}
\toprule
\textbf{ID} & \textbf{Descrição} & \textbf{Prioridade} & \textbf{Sprint} \\
\midrule
RF23 & O utilizador deve poder definir orçamentos mensais por categoria & Must & 6 \\
RF24 & O sistema deve alertar quando um orçamento atinge 80\% e 100\% & Must & 6 \\
RF25 & O utilizador deve poder criar metas de poupança com valor alvo e prazo & Must & 7 \\
RF26 & O sistema deve mostrar progresso visual das metas & Must & 7 \\
RF27 & O sistema deve sugerir montantes de poupança com base no histórico & Could & 7 \\
\bottomrule
\end{tabularx}
\end{table}

\subsection{Módulo 6 --- Simulador IRS}

\begin{table}[H]
\centering
\caption{Requisitos funcionais do módulo de simulação IRS.}
\label{tab:rf-irs}
\small
\begin{tabularx}{\textwidth}{l X c c}
\toprule
\textbf{ID} & \textbf{Descrição} & \textbf{Prioridade} & \textbf{Sprint} \\
\midrule
RF28 & O sistema deve calcular \ac{IRS} estimado com base nos rendimentos & Must & 7 \\
RF29 & O simulador deve suportar rendimentos de Categoria A & Must & 7 \\
RF30 & O simulador deve considerar deduções básicas (saúde, educação, habitação) & Must & 7 \\
RF31 & O sistema deve alertar sobre deduções que o utilizador pode estar a perder & Should & 7 \\
RF32 & O simulador deve comparar \ac{IRS} conjunto e separado (casais) & Could & 7 \\
\bottomrule
\end{tabularx}
\end{table}

\subsection{Módulo 7 --- Notificações}

\begin{table}[H]
\centering
\caption{Requisitos funcionais do módulo de notificações.}
\label{tab:rf-notif}
\small
\begin{tabularx}{\textwidth}{l X c c}
\toprule
\textbf{ID} & \textbf{Descrição} & \textbf{Prioridade} & \textbf{Sprint} \\
\midrule
RF33 & O sistema deve enviar notificações por email para alertas de orçamento & Should & 8 \\
RF34 & O sistema deve notificar sobre transações invulgares & Could & 8 \\
RF35 & O sistema deve enviar alertas sazonais (prazo IRS, IMI, IUC) & Could & 8 \\
\bottomrule
\end{tabularx}
\end{table}


\section{Requisitos Não-Funcionais}
\label{sec:req-nao-funcionais}

\subsection{Performance}

\begin{table}[H]
\centering
\caption{Requisitos não-funcionais de performance.}
\label{tab:rnf-perf}
\begin{tabularx}{\textwidth}{l X c}
\toprule
\textbf{ID} & \textbf{Descrição} & \textbf{Métrica} \\
\midrule
RNF01 & Tempo de carregamento da página inicial & $< 2$ segundos \\
RNF02 & Tempo de resposta das \ac{API}s & $< 500$ms (p95) \\
RNF03 & Suporte de utilizadores simultâneos & Mínimo 100 \\
RNF04 & Taxa de frames na \ac{UI} (scroll, animações) & $> 55$ FPS \\
\bottomrule
\end{tabularx}
\end{table}

\subsection{Segurança}

\begin{table}[H]
\centering
\caption{Requisitos não-funcionais de segurança.}
\label{tab:rnf-seg}
\begin{tabularx}{\textwidth}{l X c}
\toprule
\textbf{ID} & \textbf{Descrição} & \textbf{Métrica} \\
\midrule
RNF05 & Encriptação de dados em trânsito & \ac{TLS} 1.3 \\
RNF06 & Encriptação de dados sensíveis em repouso & AES-256 \\
RNF07 & Conformidade com \ac{RGPD} & Consentimento explícito \\
RNF08 & Proteção contra \ac{OWASP} Top 10 & Auditoria pré-release \\
RNF09 & Rate limiting nas \ac{API}s & Máx. 100 req/min \\
\bottomrule
\end{tabularx}
\end{table}

\subsection{Usabilidade}

\begin{table}[H]
\centering
\caption{Requisitos não-funcionais de usabilidade.}
\label{tab:rnf-usab}
\begin{tabularx}{\textwidth}{l X c}
\toprule
\textbf{ID} & \textbf{Descrição} & \textbf{Métrica} \\
\midrule
RNF10 & Pontuação \ac{SUS} & $> 70$ (``Good'') \\
RNF11 & Acessibilidade \ac{WCAG} & Nível AA \\
RNF12 & Suporte de browsers & Chrome, Firefox, Safari, Edge \\
RNF13 & Design responsivo & 375px a 2560px \\
RNF14 & Contraste sobre glass effects & Rácio mín. 4.5:1 \\
\bottomrule
\end{tabularx}
\end{table}

\subsection{Fiabilidade e IA}

\begin{table}[H]
\centering
\caption{Requisitos não-funcionais de fiabilidade e IA.}
\label{tab:rnf-fiab}
\begin{tabularx}{\textwidth}{l X c}
\toprule
\textbf{ID} & \textbf{Descrição} & \textbf{Métrica} \\
\midrule
RNF15 & Disponibilidade do sistema & $> 99\%$ \\
RNF16 & Cobertura de testes unitários & $> 70\%$ \\
RNF17 & Backup de dados & Diário, automático \\
RNF18 & Precisão da categorização automática & $> 80\%$ accuracy \\
RNF19 & Precisão do simulador \ac{IRS} & 10 cenários validados \\
RNF20 & Tempo de categorização por transação & $< 200$ms \\
\bottomrule
\end{tabularx}
\end{table}


\section{Diagrama de Casos de Uso}
\label{sec:req-casos-uso}

A Figura~\ref{fig:casos-uso} apresenta o diagrama de casos de uso principal do sistema, identificando os atores e as funcionalidades principais.

% NOTA: O diagrama será gerado com TikZ - ver ficheiro diagramas/casos-uso.tex
\begin{figure}[H]
\centering
% Diagrama de Casos de Uso - Gold Lock (corrigido)
\resizebox{\textwidth}{!}{%
\begin{tikzpicture}[
    usecase/.style={draw, ellipse, minimum width=3.2cm, minimum height=1.1cm, text width=2.8cm, align=center, font=\small},
    system/.style={draw, dashed, rounded corners, inner sep=20pt},
    >=Stealth
]

% Sistema boundary
\node[system, minimum width=8cm, minimum height=13cm] (sys) at (5, -3) {};
\node[font=\bfseries, anchor=south] at (5, 3.8) {Gold Lock --- Plataforma de Gestão Financeira};

% Casos de uso - bem espaçados verticalmente
\node[usecase] (uc1) at (5, 3) {Registar Conta};
\node[usecase] (uc2) at (5, 1.5) {Fazer Login};
\node[usecase] (uc3) at (5, 0) {Ligar Conta Bancária};
\node[usecase] (uc4) at (5, -1.5) {Consultar Dashboard};
\node[usecase] (uc5) at (5, -3) {Gerir Orçamentos};
\node[usecase] (uc6) at (5, -4.5) {Simular IRS};
\node[usecase] (uc7) at (5, -6) {Definir Metas Poupança};

% Ator Utilizador (à esquerda, centrado)
\draw (-1, -0.8) circle (0.35);
\draw (-1, -1.15) -- (-1, -2.1);
\draw (-1.5, -1.55) -- (-0.5, -1.55);
\draw (-1, -2.1) -- (-1.4, -2.8);
\draw (-1, -2.1) -- (-0.6, -2.8);
\node[font=\small\bfseries, anchor=north] at (-1, -2.9) {Utilizador};

% Ator Salt Edge (à direita)
\draw (11, -0.8) circle (0.35);
\draw (11, -1.15) -- (11, -2.1);
\draw (10.5, -1.55) -- (11.5, -1.55);
\draw (11, -2.1) -- (10.6, -2.8);
\draw (11, -2.1) -- (11.4, -2.8);
\node[font=\small\bfseries, anchor=north] at (11, -2.9) {Salt Edge API};

% Ligações utilizador -> casos de uso
\draw[->] (-0.5, -1.5) -- (uc1.west);
\draw[->] (-0.5, -1.5) -- (uc2.west);
\draw[->] (-0.5, -1.5) -- (uc3.west);
\draw[->] (-0.5, -1.5) -- (uc4.west);
\draw[->] (-0.5, -1.5) -- (uc5.west);
\draw[->] (-0.5, -1.5) -- (uc6.west);
\draw[->] (-0.5, -1.5) -- (uc7.west);

% Ligação Salt Edge -> Ligar Conta Bancária
\draw[->] (10.5, -1.5) -- (uc3.east);

% Relações include (com espaço para label)
\draw[->, dashed] (uc3.north) -- node[right, font=\scriptsize, fill=white, inner sep=1pt] {$\ll$include$\gg$} (uc2.south);
\draw[->, dashed] (uc4.north) -- node[right, font=\scriptsize, fill=white, inner sep=1pt] {$\ll$include$\gg$} (uc3.south);

\end{tikzpicture}%
}

\caption{Diagrama de casos de uso do sistema FinTwin.}
\label{fig:casos-uso}
\end{figure}


\section{Conclusões}
\label{sec:req-conclusoes}

A análise de requisitos resultou na identificação de 35 requisitos funcionais e 20 requisitos não-funcionais, organizados em 7 módulos funcionais. A priorização MoSCoW permite focar o desenvolvimento nos requisitos \textit{Must} durante as primeiras sprints, garantindo a entrega de um \ac{MVP} funcional, enquanto os requisitos \textit{Should} e \textit{Could} são implementados progressivamente conforme a disponibilidade de tempo.
